\section{Experimental Setup}
\label{sec:experiments}

Unless marked otherwise, the choices below are provisional.

\subsection{Dataset}
\label{sec:dataset}

Two public datasets are currently considered:
\begin{itemize}
\item \textbf{BUSI:} Breast ultrasound images labeled as normal, benign, or malignant, with mask annotations~\cite{aldhabyani2020busi}.
\item \textbf{CBIS-DDSM:} Curated mammography cases with pathology information and, for many abnormalities, BI-RADS-related descriptors such as mass shape/margin and calcification type/distribution~\cite{lee2017cbis}. The current implementation uses CBIS-DDSM.
\end{itemize}
Mass and calcification tables are merged. After that merge, some patients appear in both official train and test splits; those patients are assigned entirely to train so that evaluation remains patient-disjoint. As stated in Section~\ref{sec:concepts}, pathology is reduced to a binary label by treating benign and benign-without-callback as a single negative class. This is a simplifying assumption. Selection of any additional dataset will depend on label availability, image quality, class distribution, concept-relevant annotations, and compatibility with the proposed representation.
% TODO: Select the primary dataset; report splits, class counts, preprocessing, and exclusion criteria.

\begin{table}[!t]
\caption{Dataset characteristics (to be completed after final selection).}
\label{tab:dataset}
\centering
\small
\begin{tabular}{@{}lccc@{}}
\toprule
Dataset & Modality & Samples & Concept labels \\
\midrule
BUSI~\cite{aldhabyani2020busi} & Ultrasound & \todo{N} & \todo{verify} \\
CBIS-DDSM~\cite{lee2017cbis} & Mammography & \todo{N} & \todo{verify} \\
\bottomrule
\end{tabular}
\end{table}

\subsection{Neural Model}
\label{sec:neural}

Candidate backbones include a standard CNN baseline and residual networks such as ResNet~\cite{he2016resnet}; alternatives may be considered. Geometric augmentation (horizontal flip, small rotation, and mild affine perturbation) is applied to training images only; validation and test images are left unchanged. Learning rate and batch size are selected by a small grid search on validation macro F1, after which the selected pair is used for the full training run.
Two neural roles are planned: an end-to-end diagnostic classifier (baseline) and a concept extractor for the neuro-symbolic pipeline.

\subsection{DeepProbLog Configuration}
\label{sec:dplconfig}

The planned configuration includes neural predicates for lesion kind, shape, margin, and calcification type; probabilistic facts from those network outputs; symbolic rules that gate mass evidence and calcification evidence by lesion kind; and a diagnostic query for malignancy probability. These rules are a computational working program and are not presented as official BI-RADS criteria.
% TODO: Specify the final program, predicate signatures, query, learning schedule, and inference settings.

\subsection{Baselines}
\label{sec:baselines}

\begin{itemize}
\item \textbf{Baseline 1:} Pure neural image classifier.
\item \textbf{Baseline 2:} Neural concept-level model without DeepProbLog, if annotations permit (e.g., concept-bottleneck style)~\cite{koh2020concept}.
\item \textbf{Proposed:} Neural concepts + DeepProbLog + symbolic medical knowledge.
\end{itemize}

\subsection{Evaluation Metrics}
\label{sec:metrics}

The evaluation will report accuracy, precision, recall/sensitivity, specificity, F1-score for both classes (and their macro average), and ROC-AUC where applicable, with emphasis on sensitivity and specificity. Reporting only malignant F1 would hide a model that predicts a single class. Interpretability will be assessed qualitatively via reasoning traces on selected true/false positive and negative cases; expert review may be added if resources permit.
% TODO: Define the interpretability annotation protocol.
